% Copie os arquivos PDF para a pasta images/matlab do seu relatorio.
% Figuras geradas a partir de simulacao numerica no MATLAB.
\begin{figure}[htbp]
\centering
\includegraphics[width=\linewidth]{images/matlab/01_Fig1_sinal1.pdf}
\caption{Caso 01: Rampa crescente unipolar. Comparacao entre teoria e simulacao: onda, FFT e energias.}
\label{fig:matlab_caso_01}
\end{figure}

\begin{figure}[htbp]
\centering
\includegraphics[width=\linewidth]{images/matlab/02_Fig1_sinal2.pdf}
\caption{Caso 02: Rampa decrescente unipolar. Comparacao entre teoria e simulacao: onda, FFT e energias.}
\label{fig:matlab_caso_02}
\end{figure}

\begin{figure}[htbp]
\centering
\includegraphics[width=\linewidth]{images/matlab/03_Fig1_sinal3.pdf}
\caption{Caso 03: Triangular unipolar. Comparacao entre teoria e simulacao: onda, FFT e energias.}
\label{fig:matlab_caso_03}
\end{figure}

\begin{figure}[htbp]
\centering
\includegraphics[width=\linewidth]{images/matlab/04_Fig1_sinal4.pdf}
\caption{Caso 04: Rampa crescente bipolar. Comparacao entre teoria e simulacao: onda, FFT e energias.}
\label{fig:matlab_caso_04}
\end{figure}

\begin{figure}[htbp]
\centering
\includegraphics[width=\linewidth]{images/matlab/05_Fig1_sinal5.pdf}
\caption{Caso 05: Rampa decrescente bipolar. Comparacao entre teoria e simulacao: onda, FFT e energias.}
\label{fig:matlab_caso_05}
\end{figure}

\begin{figure}[htbp]
\centering
\includegraphics[width=\linewidth]{images/matlab/06_Fig1_sinal6.pdf}
\caption{Caso 06: Triangular bipolar. Comparacao entre teoria e simulacao: onda, FFT e energias.}
\label{fig:matlab_caso_06}
\end{figure}

\begin{figure}[htbp]
\centering
\includegraphics[width=\linewidth]{images/matlab/07_Fig2_1_0P_sinal1.pdf}
\caption{Caso 07: Rampa crescente unipolar - $\delta=P$, $\phi=P/10$. Comparacao entre teoria e simulacao: onda, FFT e energias.}
\label{fig:matlab_caso_07}
\end{figure}

\begin{figure}[htbp]
\centering
\includegraphics[width=\linewidth]{images/matlab/08_Fig2_1_0P_sinal2.pdf}
\caption{Caso 08: Rampa decrescente unipolar - $\delta=P$, $\phi=P/10$. Comparacao entre teoria e simulacao: onda, FFT e energias.}
\label{fig:matlab_caso_08}
\end{figure}

\begin{figure}[htbp]
\centering
\includegraphics[width=\linewidth]{images/matlab/09_Fig2_1_0P_sinal3.pdf}
\caption{Caso 09: Triangular unipolar - $\delta=P$, $\phi=P/10$. Comparacao entre teoria e simulacao: onda, FFT e energias.}
\label{fig:matlab_caso_09}
\end{figure}

\begin{figure}[htbp]
\centering
\includegraphics[width=\linewidth]{images/matlab/10_Fig2_1_0P_sinal4.pdf}
\caption{Caso 10: Rampa crescente bipolar - $\delta=P$, $\phi=P/10$. Comparacao entre teoria e simulacao: onda, FFT e energias.}
\label{fig:matlab_caso_10}
\end{figure}

\begin{figure}[htbp]
\centering
\includegraphics[width=\linewidth]{images/matlab/11_Fig2_1_0P_sinal5.pdf}
\caption{Caso 11: Rampa decrescente bipolar - $\delta=P$, $\phi=P/10$. Comparacao entre teoria e simulacao: onda, FFT e energias.}
\label{fig:matlab_caso_11}
\end{figure}

\begin{figure}[htbp]
\centering
\includegraphics[width=\linewidth]{images/matlab/12_Fig2_1_0P_sinal6.pdf}
\caption{Caso 12: Triangular bipolar - $\delta=P$, $\phi=P/10$. Comparacao entre teoria e simulacao: onda, FFT e energias.}
\label{fig:matlab_caso_12}
\end{figure}

\begin{figure}[htbp]
\centering
\includegraphics[width=\linewidth]{images/matlab/13_Fig2_0_8P_sinal1.pdf}
\caption{Caso 13: Rampa crescente unipolar - $\delta=0{,}8P$, $\phi=P/10$. Comparacao entre teoria e simulacao: onda, FFT e energias.}
\label{fig:matlab_caso_13}
\end{figure}

\begin{figure}[htbp]
\centering
\includegraphics[width=\linewidth]{images/matlab/14_Fig2_0_8P_sinal2.pdf}
\caption{Caso 14: Rampa decrescente unipolar - $\delta=0{,}8P$, $\phi=P/10$. Comparacao entre teoria e simulacao: onda, FFT e energias.}
\label{fig:matlab_caso_14}
\end{figure}

\begin{figure}[htbp]
\centering
\includegraphics[width=\linewidth]{images/matlab/15_Fig2_0_8P_sinal3.pdf}
\caption{Caso 15: Triangular unipolar - $\delta=0{,}8P$, $\phi=P/10$. Comparacao entre teoria e simulacao: onda, FFT e energias.}
\label{fig:matlab_caso_15}
\end{figure}

\begin{figure}[htbp]
\centering
\includegraphics[width=\linewidth]{images/matlab/16_Fig2_0_8P_sinal4.pdf}
\caption{Caso 16: Rampa crescente bipolar - $\delta=0{,}8P$, $\phi=P/10$. Comparacao entre teoria e simulacao: onda, FFT e energias.}
\label{fig:matlab_caso_16}
\end{figure}

\begin{figure}[htbp]
\centering
\includegraphics[width=\linewidth]{images/matlab/17_Fig2_0_8P_sinal5.pdf}
\caption{Caso 17: Rampa decrescente bipolar - $\delta=0{,}8P$, $\phi=P/10$. Comparacao entre teoria e simulacao: onda, FFT e energias.}
\label{fig:matlab_caso_17}
\end{figure}

\begin{figure}[htbp]
\centering
\includegraphics[width=\linewidth]{images/matlab/18_Fig2_0_8P_sinal6.pdf}
\caption{Caso 18: Triangular bipolar - $\delta=0{,}8P$, $\phi=P/10$. Comparacao entre teoria e simulacao: onda, FFT e energias.}
\label{fig:matlab_caso_18}
\end{figure}

\begin{figure}[htbp]
\centering
\includegraphics[width=\linewidth]{images/matlab/19_Ret7_simetrica.pdf}
\caption{Caso 19: Retangular unipolar simetrica. Comparacao entre teoria e simulacao: onda, FFT e energias.}
\label{fig:matlab_caso_19}
\end{figure}

\begin{figure}[htbp]
\centering
\includegraphics[width=\linewidth]{images/matlab/20_Ret7_assimetrica.pdf}
\caption{Caso 20: Retangular unipolar assimetrica. Comparacao entre teoria e simulacao: onda, FFT e energias.}
\label{fig:matlab_caso_20}
\end{figure}

\begin{figure}[htbp]
\centering
\includegraphics[width=\linewidth]{images/matlab/21_Ret7_alternada.pdf}
\caption{Caso 21: Retangular unipolar alternada. Comparacao entre teoria e simulacao: onda, FFT e energias.}
\label{fig:matlab_caso_21}
\end{figure}

\begin{figure}[htbp]
\centering
\includegraphics[width=\linewidth]{images/matlab/22_Ret8_simetrica.pdf}
\caption{Caso 22: Retangular bipolar simetrica. Comparacao entre teoria e simulacao: onda, FFT e energias.}
\label{fig:matlab_caso_22}
\end{figure}

\begin{figure}[htbp]
\centering
\includegraphics[width=\linewidth]{images/matlab/23_Ret8_assimetrica.pdf}
\caption{Caso 23: Retangular bipolar assimetrica. Comparacao entre teoria e simulacao: onda, FFT e energias.}
\label{fig:matlab_caso_23}
\end{figure}

\begin{figure}[htbp]
\centering
\includegraphics[width=\linewidth]{images/matlab/24_Ret8_alternada.pdf}
\caption{Caso 24: Retangular bipolar alternada. Comparacao entre teoria e simulacao: onda, FFT e energias.}
\label{fig:matlab_caso_24}
\end{figure}

